\appendixsection{Operations on linear representations}


Let \(G\) be a group, \(k\) a commutative field, \(E_1, E_2, \ldots, E_n\) vector spaces over \(k\) and \(\pi_i\) a linear representation of \(G\) on \(E_i\) ( \(1 \leqslant i \leqslant n\) ). The mapping

\[
g \mapsto \pi_ {1} (g) \otimes \dots \otimes \pi_ {n} (g)
\]

is a linear mapping of G into the vector space \(E_{1} \otimes \cdots \otimes E_{n}\) , called the tensor product of \(\pi_{1}, \ldots, \pi_{n}\) and denoted by \(\pi_{1} \otimes \cdots \otimes \pi_{n}\) .

Let E be a vector space over k and \(\pi\) a linear representation of G on E. For all \(g \in G\) , let \(\tau(g)\) (resp. \(\sigma(g)\) , \(\varepsilon(g)\) ) be the unique automorphism of the algebra \(\mathsf{T}(\mathsf{E})\) (resp. \(\mathsf{S}(\mathsf{E})\) , \(\bigwedge(\mathsf{E})\) ) which extends \(\pi(g)\) (Algebra, Chapter III, § 5, no. 2, § 6, no. 2 and § 7, no. 2). Then \(\tau\) (resp. \(\sigma\) , \(\varepsilon\) ) is a linear representation of G on T(E) (resp. \(\mathsf{S}(\mathsf{E})\) , \(\bigwedge(\mathsf{E})\) ) denoted by \(\mathsf{T}(\pi)\) (resp. \(\mathsf{S}(\pi)\) , \(\bigwedge(\pi)\) ). The subrepresentation of T(π) (resp. \(\mathsf{S}(\pi)\) , \(\bigwedge(\pi)\) ) defined by \(\mathsf{T}^{\mathsf{n}}(\mathsf{E})\) (resp. \(\mathsf{S}^{\mathsf{n}}(\mathsf{E})\) , \(\bigwedge^{\mathsf{n}}(\mathsf{E})\) ) is called the n-th tensor (resp. symmetric, exterior) power of \(\pi\) and is denoted by \(\mathsf{T}^{\mathsf{n}}(\pi)\) (resp. \(\mathsf{S}^{\mathsf{n}}(\pi)\) , \(\bigwedge^{\mathsf{n}}(\pi)\) ). Then

\[
T ^ {n} (\pi) = \pi \otimes \pi \otimes \dots \otimes \pi
\]

(n factors). The representations \(\mathsf{S}(\pi), \bigwedge (\pi)\) are quotient representations of \(\mathsf{T}(\pi)\) and hence \(\mathsf{S}^n (\pi), \bigwedge^n (\pi)\) are quotient representations of \(\mathsf{T}^n (\pi)\) .

Let g be a Lie algebra over k. The tensor product of a finite number of representations of g has already been defined in Chapter I, §3, no. 2; it is denoted by \(\pi_{1}\otimes\cdots\otimes\pi_{n}\) . Let E be a vector space over k and \(\pi\) a representation of g on E. For all \(x\in g\) , let \(\tau'(x)\) (resp. \(\sigma'(x)\) , \(\varepsilon'(x)\) ) be the unique derivation of the algebra \(\mathrm{T(E)}\) (resp. \(\mathrm{S(E)}\) , \(\bigwedge(\mathrm{E})\) ) which extends \(\pi(x)\) (Algebra, Chapter III, §10, no. 9, Example 1). Then \(\tau'\) (resp. \(\sigma'\) , \(\varepsilon'\) ) is a linear representation of g on T(E) (resp. S(E), \(\bigwedge(\mathrm{E})\) ) by Algebra, loc. cit., formula (35),

denoted by \(\mathsf{T}(\pi)\) (resp. \(\mathsf{S}(\pi),\bigwedge (\pi)\) ). The subrepresentation of \(\mathsf{T}(\pi)\) (resp. \(\mathsf{S}(\pi),\bigwedge (\pi)\) ) defined by \(\mathsf{T}^n (\mathsf{E})\) (resp. \(\mathsf{S}^n (\mathsf{E}),\bigwedge^n (\mathsf{E}))\) is denoted by \(\mathsf{T}^n (\pi)\) (resp. \(\mathsf{S}^n (\pi),\bigwedge^n (\pi))\) . The representation \(\mathsf{T}^n (\pi)\) is the tensor product on \(n\) representations identical with \(\pi\) . The representations \(\mathsf{S}(\pi),\bigwedge (\pi)\) are quotient representations of \(\mathsf{T}(\pi)\) and hence \(\mathsf{S}^n (\pi),\bigwedge^n (\pi)\) are quotient representations of \(\mathsf{T}^n (\pi)\) .

